\section{5교시 (75분) — 품질 · 이슈 · 리스크 · 조달 통제}

% =====================================================================
% 5교시 — 품질 · 이슈 · 리스크 · 조달 통제 (교재 제4장·제5장, 워크북 §3~§5)
% =====================================================================
\begin{frame}{5교시에서 배우는 것 — 나머지 세 산출물(⑬·⑭·⑮)}
\begin{itemize}
\item 품질 — 결함밀도의 분모·분자를 누가 정했는가, 곡선의 평탄은 수렴인가 정지인가
\item 감리 — 등급보다 \textbf{원인 귀속} 칸이 나중의 책임 귀속을 정한다
\item 이슈·리스크 — 이슈에는 기한, 리스크에는 트리거. 잔여 EMV 7.26 vs 허용 8.0
\item 보고 — RAG의 정치와 mum effect. 나쁜 소식은 왜 올라오지 않는가
\item 계약 — 동시기록(검수 보류 서면 회신)·동시 지연 분해·수행사 성과 기록
\end{itemize}
\keymsg{네 번째 산출물부터는 숫자보다 "누가 원인이고 누가 결정하는가"를 적는 칸이 무겁다}
\note{교재 제4장(품질·이슈·리스크)과 제5장(사람·커뮤니케이션·조달)을 75분에 다룬다. 워크북 §3(⑬ 이슈 로그·리스크 등록부 갱신), §4(⑭ 품질 통제 기록), §5(⑮ 조달·계약 변경 기록)의 완성 예시본이 정본이다. 1\til{}4교시가 "편차를 재는 법"이었다면 5교시는 "편차의 원인을 누구에게 귀속시키고, 그 기록이 검수·변경협의체·지체상금에서 어떻게 쓰이는가"다. 세 산출물은 서로를 인용한다 — 감리 지적 7번(⑭)이 이슈 IS-07(⑬)이고 R-20(⑬)이며 검수 보류 회신(⑮)의 4항이다. 수강생에게 워크북 §4.5 항목 6 대응 관리대장과 §5.5 항목 4 서면 회신을 펴 두게 한다. 예상 소요 4분.}
\end{frame}

\begin{frame}{QA와 QC, 그리고 결함밀도 정의의 정치 — 분모와 분자}
\begin{itemize}
\item QA는 절차가 지켜졌는가(감리·검토회), QC는 산출물이 기준을 넘었는가(결함·밀도)
\item 분자 — 스프린트 시험 결함만(설계 검토 제외), 요구 변경·해석 차이는 결함이 아니라 CR
\item 분모 — \textbf{인수 FP 누적} 7,020. 계획 FP(7,440)로 나누면 0.125로 "이내"가 된다
\end{itemize}
\begin{tbl}[\scriptsize]{L{3.2cm}L{3.6cm}rL{2.3cm}L{2.6cm}}
계층 & 산식 & 값 & 기준 & 판정\\\midrule
작업패키지 & 930 $\div$ 7,020 FP & 0.132 & 0.13 (수행사 PM) & \textbf{초과 +1.5\%} \ra{} 1단계\\
설계 검토 포함(참고) & 1,140 $\div$ 7,020 & 0.162 & — & 분모 부정합\\
프로젝트(전망) & 0.132 $\times$ 유출 35\% + UAT 0.04 & 약 0.09 & 0.4 + 심1·2 미해결 0 & 이내 전망\\
\end{tbl}
\keymsg{0.132 > 0.13 — 1.5\% 초과라도 허용범위 밖이면 에스컬레이션한다. 반올림으로 "이내"를 만들지 않는다}
\note{교재 4.1, 워크북 §4.5 항목 1·4. 결함밀도는 정의가 곧 정치다 — 분자에 설계 검토 결함을 넣으면 0.162, 분모를 계획 FP로 바꾸면 0.125가 되고, 세 숫자 중 어느 것이 "기준 0.13"과 비교되는지는 Day2 격자에서 정했다(스프린트 시험 결함 ÷ 인수 FP). 그래서 문서 항목 1에 측정 정의를 먼저 적는다. 0.132는 0.13을 1.5\% 넘었을 뿐이지만 예외에 의한 관리에서 허용범위는 허용범위다 — 이재현이 인지 09-25 다음 영업일 09-28에 P1 PM에 보고했고, P1 PM은 3영업일 내(10-01)에 결함 해소 전담 2명 배치·backlog 10-31 60 이하·시험 노력 주간 보고 추가를 결정했다. 프로젝트 계층 0.4는 넉넉하고 실질 제약은 "심각도 1·2 미해결 0"임을 짚는다. 예상 소요 6분.}
\end{frame}

\begin{frame}{결함 누적곡선 — 평탄에는 두 가지 뜻이 있다}
\fig[0.58]{../그림/PM_Day3/fig_4_2_defect_curves.pdf}
\figcap{그림 4-2. 34주 누적 발견 1,140 · 해소 1,024 · 미해결 116. 3월 말 평탄(①)은 수렴, 8월 평탄(③)은 주당 TC 210 \ra{} 95로 시험이 멈춘 것 — 곡선만 보면 둘은 같아 보인다}
\keymsg{발견 곡선이 평평해지면 먼저 시험 노력을 본다. 노력 없이 줄어든 결함은 품질이 아니라 정지다}
\note{교재 4.2, 워크북 §4.5 항목 2·3. 네 구간을 읽는다 — ① 설계 검토 2\til{}3월 주 18\ra{}34\ra{}15, 검토 대상 47종이 S5까지 대부분 끝나 노력 감소 없이 발견이 줄었으므로 수렴형. ② S7\til{}S15 인수 FP에 비례해 상승, 6월 밀도 0.16 최고(R-05 지연 후 S11·S12 재작업). ③ 8월 12\ra{}9\ra{}11 — 시험 인력 6명 중 4명이 규제 검증(A07 142개 항목)으로 이동해 주당 TC가 210에서 95로 떨어졌다. 8월 밀도 0.129는 착시다. ④ 9월 재상승 27\ra{}52\ra{}28, 8월 미시험분이 S17에 몰렸다. 지금 봐야 할 것은 발견이 아니라 backlog다 — 6월부터 해소가 발견을 못 따라가 116건이 쌓였고 R3 개발자 시간을 잠식한다(A10 $-$10의 원인 중 하나). 10월부터 시험 노력(TC 실행)을 결함표 열로 추가하는 이유다. 예상 소요 6분.}
\end{frame}

\begin{frame}{심각도 2 미해결 7건 — 건수보다 나이를 본다}
\begin{tbl}[\scriptsize]{lL{3.6cm}rlL{3.4cm}}
ID & 모듈 (WBS) & 나이 & 기한 & 연결\\\midrule
D-0966 & API GW 인증 (1.5.2) & 16 & 10-16 & IS-08 보안 High, 감리 6번\\
D-0871 & 3PL 폴백 전환 (1.3.3) & 12 & 10-09 & R-16, TC-INV-037\\
D-0902 & POS 반품 적립 취소 (1.1.2) & 10 & 10-09 & 멤버십 StRS\\
D-0915 & 개인정보 파기 로그 (1.5.1) & 8 & 10-16 & 감리 6번, SRS-708\\
D-0958 & MDM 중복 검출 (1.5.3) & 4 & 10-23 & R-10, 정본 판정 전 필수\\
\end{tbl}
\begin{itemize}
\item 미해결 116 = 심1 \textbf{0} / 심2 \textbf{7} / 심3 58 / 심4 51 — 인수에는 지장 없으나 시간을 잠식
\item 나이(영업일) 중앙값 9, 심2 최장 16 — "심2 나이 10일 초과"를 주간 PM 회의 고정 안건으로
\item 심각도 기준서 v1.0(09-28)은 감리 5번 대응 — 116건 재분류 10-09
\end{itemize}
\keymsg{통합시험 착수(12-07) 조건은 심2 미해결 0·심3 30 이하 — G3의 "심1·2 = 0"은 지금부터 만든다}
\note{교재 4.2 "심각도와 나이", 워크북 §4.5 항목 5(7건 중 5건만 표시 — D-0931·D-0944는 워크북 참조). 결함 건수는 인수 FP에 비례해 늘기 마련이므로 절대 건수는 정보가 적다. 정보가 있는 것은 심각도별 미해결 수와 그 나이다 — D-0966은 보안 High(VA-1-01)와 같은 건이고 16영업일째 열려 있으며, D-0915는 감리 6번(파기 증적)과 원인이 같다. 결함 하나가 이슈(IS-08)·감리 지적·리스크(R-16·R-10)에 동시에 연결된다는 점을 보여 준다 — 그래서 연결 열이 있다. 우선순위(P1\til{}P3)는 심각도와 별도 열이라는 것도 짚는다. 예상 소요 5분.}
\end{frame}

\begin{frame}{감리 지적 8건 — 등급보다 원인 귀속 칸이 무겁다}
\fig[0.56]{../그림/PM_Day3/fig_4_3_audit_attribution.pdf}
\figcap{그림 4-3. 감리 기재 수행사 6·발주사 2 \ra{} 발주사 판정 수행사 3·발주사 4·공동 1. AF-1-03(P2 규격 확인 서명)·AF-1-04(가맹 SW 배포 범위)는 WBS 담당 열·RACI로 발주사 귀속을 입증해 09-29 서면 이의}
\keymsg{지적 사실은 인정하고 귀속만 다툰다 — 그 두 글자가 R2 검수·대기 비용·변경협의체의 근거가 된다}
\note{교재 4.3, 워크북 §4.5 항목 6과 이의 제기 서면 완성본. 감리 지적 대응 관리대장의 핵심 열은 등급(중대·보통·경미)이 아니라 "감리 기재 귀속"과 "발주사 판정 귀속"의 두 열이다. AF-1-03은 규격이 WBS 1.3.1 발주 담당 산출물이고 P2 전달은 RACI 9행 A = 프로그램 PMO장인데, 설비 벤더의 05-29 변경 요청에 프로그램 PMO가 09-30 현재 회신하지 않았다(회신 요청 06-05·07-10·08-21, 수행사 미결 통지 06-08·08-24 기록). AF-1-04는 1.1.6 대응안 결정이 오정근 A이며 05-08 "부속서 협상 후로 유보" 회신이 있다. 이의 서면의 구조 — ① 지적 사실 인정 ② 귀속 재분류 요청과 근거(WBS 담당 열·RACI·기록) ③ 발주사 조치 계획(10-16·11-27) ④ 요청(정본 정정, 10-08 회신, 향후 귀속 판정 기준 합의). 등급을 다투지 않는 이유 — 등급은 감리의 권한이고, 귀속은 사실의 문제다. 예상 소요 6분.}
\end{frame}

\begin{frame}{이슈 로그 11건 — 이슈에는 기한이, 리스크에는 트리거가 있다}
\begin{tbl}[\scriptsize]{lL{4.2cm}lL{2.2cm}l}
ID & 제목 & 심 & 에스컬 & 상태\\\midrule
IS-01 & R-04 실현 — 2026 배분 $-$8.0 & 1 & 2단계(스폰서) & 종결 03-13\\
IS-02 & R-05 실현 — POS 화면 승인 12영업일 & 1 & 병행(변경협의체) & 종결 07-15\\
IS-05 & R-03 트리거 — 3PL 폴백 확정 & 2 & 병행(임원\ra{}스폰서) & 종결 08-21\\
IS-07 & 감리 8건 — 3·4번 귀속 이의 & 1 & 2단계(회신 10-08) & 진행\\
IS-08 & 보안진단 High 2건 & 1 & 병행(최영주) & 진행\\
IS-11 & 수행사 인터페이스 PL 교체 & 2 & 1단계 & 종결 08-03\\
\end{tbl}
\begin{itemize}
\item 출처 — 리스크 전환 3 / 신규 발견 2 / 감리·진단 2 / 외부 3 / 수행사 1 (종결 5 · 진행 6)
\item 등록부에 없던 원인 5건 중 3건(IS-04·08·11)은 프리모템 원문에 있었다 \ra{} Day4 교훈 ⑲
\item 트리거 발동은 신호(기록·집행 준비), 실현은 사건(행 이동·컨틴전시 집행·이슈 개설)
\end{itemize}
\keymsg{에스컬레이션 단계마다 응답 시한이 있다 — 1단계 3영업일 · 2단계 10영업일. IS-05의 8월 CCB 보고 누락이 감리 8번이 됐다}
\note{교재 4.4, 워크북 §3.5 항목 2·3·8(표는 11건 중 6건 발췌). 이슈 로그의 열 규칙 — 발생일과 인지일을 따로(IS-02는 05-12 발생, 05-14 인지), 소유자와 조치자를 따로, 기한과 에스컬레이션 단계를 따로. 리스크 \ra{} 이슈 전환 규칙: 트리거 발동은 신호이므로 소유자가 관측 사실과 일자를 등록부에 적고 집행 준비를 하지만 컨틴전시를 쓰지 않는다. 실현은 확률 100\%가 된 사건이므로 행을 "실현"으로 옮기고 이슈를 개설하며 컨틴전시를 집행한다 — 이때 CCB 보고(RACI 19행 I)가 절차인데 IS-05 R-03 0.50 집행이 8월 CCB에서 빠졌고 그것이 감리 AF-1-08이 되어 10-15 소급 보고한다. 절차 하나가 빠지면 다른 산출물에서 지적이 된다는 예다. 예상 소요 6분.}
\end{frame}

\begin{frame}{리스크 등록부 v1.0 \ra{} v1.2 — 대응은 노출을 없애지 않고 옮긴다}
\fig[0.56]{../그림/PM_Day3/fig_4_4_risk_issue_transitions.pdf}
\figcap{그림 4-4. 종결·실현·이관 4건(R-03·04·09·12), 상승 R-10(0.38 \ra{} 0.76), 신규 7건 R-14\til{}R-20 — 그중 5건(R-14·15·16·17·18)은 대응 조치가 만든 2차 리스크}
\keymsg{R-01 분리 \ra{} R-14 두 번 컷오버, R-04 이연 \ra{} R-17 2027 자금 봉우리 — 관리예비비 대상 2.20은 숫자만 같고 내용이 바뀌었다}
\note{교재 4.5 "2차 리스크의 승격", 워크북 §3.5 항목 4·5. 재평가는 세 숫자를 다시 적는 일이다 — 확률·영향·EMV. R-10(G2 후 마스터 변경)은 9월에 신규 브랜드 상품 코드·B2B 거래처 코드 변경 요구가 접수되어 10\ra{}20\%로 상승했고 컨틴전시를 0.40에서 0.80으로 증액했다. 신규 7건은 출처가 모두 기록에 있다 — R-14는 G1 R-01 분리 의결의 결과(프로그램 PMO장 소유), R-16은 R-03 폴백 배치와 CR-09 22:00 마감의 충돌, R-19는 EX-01 회복 계획(통합시험 압축 \ra{} 결함 유출), R-20은 감리 귀속 분쟁(IS-07). R-17은 원인이 스폰서 자금 결정이므로 컨틴전시가 아니라 관리예비비 대상이다. Day2에서 R-14\til{}R-18을 예비 번호로 남겨 둔 이유가 여기서 드러난다. 예상 소요 5분.}
\end{frame}

\begin{frame}{잔여 EMV 7.26 vs 8.0 — 이내지만 경고, 그리고 컨틴전시 소진 13.8\%}
\fig[0.52]{../그림/PM_Day3/fig_4_5_emv_contingency.pdf}
\figcap{그림 4-5. 6.78 $-$ 종결 1.98 + 신규 3.13 $-$ 재평가 0.67 = 7.26(여유 0.74). 컨틴전시 10.20 = 배정 5.30 + 미배정 2.47 + 집행 1.41 + BL-02 편입 1.02}
\keymsg{여유 0.74는 R-10 하나(0.76)의 크기 — 11월 갱신에서 R-10이 실현되면 즉시 8억 초과. 예외 보고 불요, 스폰서 사전 예고}
\note{교재 4.5 "컨틴전시 소진 추적", 워크북 §3.5 항목 6·7·8. 판정 규칙은 둘 — 잔여 위협 EMV 7.26 < 8.0, 단일 최대 0.80 < 2.0이므로 리스크 축 예외 보고는 불요. 그러나 여유가 1.22에서 0.74로 줄었고 R-10의 트리거(정본 판정 후 변경 요청)는 G2 직후 관측되므로 10-05 성과보고 제9호에 사전 예고한다. 컨틴전시 소진율 13.8\%(약정 1.41/10.20) vs 원가 진척 50.4\%(EV/BAC)는 얼핏 "여유롭다"로 읽히지만, 배정 재조정 후 미배정 풀이 스프린트 4회분에서 1.75회분(2.47)으로 줄었다 — 다음 신규 리스크 2건이면 풀이 마른다. 집행 1.41 중 AC 발생은 1.19(CR-07 0.22는 약정)라는 것, BL-02로 편입된 1.02는 더 이상 컨틴전시가 아니라는 것도 짚는다. 일정 축은 임계경로 잔여 영향 25.5영업일 vs 버퍼 15\til{}30으로 경고다. 예상 소요 6분.}
\end{frame}

\begin{frame}{성과보고 제9호 — RAG 일곱 축과 결정 요청 셋, 그리고 나쁜 소식의 경제학}
\fig[0.44]{../그림/PM_Day3/fig_5_2_mum_effect.pdf}
\figcap{그림 5-2. 일정 R · 원가 A · 범위 G · 품질 A · 리스크 A · 편익 G · 지속가능성 A — 결정 요청: EX-01(10-16) · CCB(10-15) · 감리 귀속 회신(10-08)}
\keymsg{RAG는 색이 아니라 문장이다 — "R"은 "결정이 필요하다"이지 "PM이 못했다"가 아니다. 그 등식이 깨진 조직에서 mum effect가 시작된다}
\note{교재 5.1\til{}5.2, 워크북 §2.5 항목 12(Highlight Report 1장). 일곱 축의 RAG를 학생이 먼저 매기게 하고(4교시 산출물 갤러리워크에서 이미 봤다) 예시본과 비교한다 — 일정 R은 IEAC(t) > 허용이므로 논쟁이 없지만 원가 A/G, 품질 A/R는 조마다 갈린다. 그 갈림이 "판정 기준을 문서에 적었는가"의 문제라는 것이 요점(워크북 §2.3 항목 12 판정 규칙). mum effect(Keil·Smith 계열 연구): 나쁜 소식은 보고자가 책임 귀속을 두려워할수록 늦어지고, 늦어질수록 비싸진다 — 온담에서는 IS-02(설계 승인 대기 12영업일)가 그 사례. 대응은 "보고자의 안전"이 아니라 "판정 기준의 사전 합의"다. 예상 소요 6분.}
\end{frame}

\begin{frame}{동시기록 — 구두 검수 보류 통보에 서면으로 답하지 않으면 발주자 사유 기간을 공제받지 못한다}
\fig[0.44]{../그림/PM_Day3/fig_5_5_contemporaneous_records.pdf}
\figcap{그림 5-5. 구두 통보(09-24) \ra{} 서면 회신(09-26, 발주사 사유 8영업일·R2 인수 FP 12\% 영향 명시) \ra{} 검수 기한 10-16 재확인 \ra{} 회의록 서명}
\keymsg{회신은 논쟁이 아니라 기록이다 — 사실(무엇을, 언제, 누가)과 영향(며칠, 몇 FP)만 적고 책임은 적지 않는다}
\note{교재 5.4 "동시기록 문화", 워크북 §5.5 항목 3(검수 보류 서면 회신 완성본). 세 가지를 짚는다. (1) 구두 통보 자체는 유효하지만 그 내용과 시점의 기록은 회신에서만 남는다 — 회신이 없으면 3개월 뒤 "보류 통보가 있었는가"부터 다툰다. (2) 회신에 반드시 들어갈 다섯 항목: 통보 사실·시점·사유의 귀속(발주사 승인 대기 8영업일)·영향(R2 인수 FP의 12\%)·재확인 요청(검수 기한). (3) 책임을 적지 않는다 — "귀사 귀책으로"가 아니라 "귀사 승인 대기 기간 8영업일"이라고 쓴다. 이것이 SCL 동시지연 분해(다음 프레임)의 재료가 되고, Day1 회의록 서명 회피(권한 상실 신호)와 같은 계열. 6교시 실습 ⑮의 과제. 예상 소요 5분.}
\end{frame}

\begin{frame}{동시지연 분해 — 발주 7 · 동시 3 · 수행 3 · 회복 $-$3, 그리고 지체상금의 계산 순서}
\fig[0.44]{../그림/PM_Day3/fig_5_6_concurrent_delay.pdf}
\figcap{그림 5-6. A10 착수 지연 10영업일의 귀책 분해 — 발주사(설계 승인 대기 IS-02) 7 · 동시(양측) 3 · 수행사(R2 이월) 3 · 수행사 회복 조치 $-$3 = 순 10}
\keymsg{지체상금은 "며칠 늦었나"가 아니라 "며칠 중 며칠이 누구 것인가"에서 시작한다 — 분해가 없으면 계약은 총일수로 청구한다}
\note{교재 5.5 "동시지연", 워크북 §5.5 항목 6. SCL Delay and Disruption Protocol 2판의 동시지연 개념을 온담 숫자로: A10 $-$10 중 발주사 사유(IS-02 설계 승인 대기 12영업일 중 임계경로에 걸린 7)·동시 3(양측 원인 겹침)·수행사 3(R2 이월 420FP)·수행사 회복 $-$3(S18 3주 연장). 지체상금 조항(일 0.05\%, 상한 10\%)은 Day1 계약 구조에서 왔고, 발주사 사유 기간은 공제되므로 순 수행사 귀책은 3 $-$ 3 = 0 — 단 이것은 서면 회신(앞 프레임)이 있을 때만 성립한다. 여유 소유권 문안(Day2 §3, "총여유는 프로젝트의 것")이 여기서 실제로 작동한다. 학생 질문 예상: "동시 3은 누가 정하나" — 답: 계약 문안 + 동시기록, 없으면 감정인. 예상 소요 5분.}
\end{frame}

\begin{frame}{수행사 성과 · 인력 교체 · 하도급 — 계약은 관계가 아니라 기록으로 관리된다}
\fig[0.44]{../그림/PM_Day3/fig_5_4_supplier_performance.pdf}
\figcap{그림 5-4. 스프린트별 velocity(S7\til{}S18)와 결함 유출, SLA 양방향(발주사 승인 SLA 포함), 인력 교체 1건(07-20 인터페이스 PL) 승인, 하도급 재위탁 4 \ra{} 5}
\keymsg{성과 기록이 양방향이어야 협상이 된다 — 수행사 velocity만 적고 발주사 승인 SLA를 안 적으면 IS-02는 수행사 탓이 된다}
\note{교재 5.3, 워크북 §5.5 항목 2·4·5. 세 기록의 요점. (1) 수행사 성과: velocity와 결함 유출률을 스프린트마다 기록하되, 발주사 측 SLA(승인 대기일·자료 제공일)를 같은 표에 둔다 — 그래야 S10\til{}S11 하락의 원인이 IS-02(발주사)임이 표에서 읽힌다. (2) 인력 교체: 07-20 인터페이스 PL 교체 승인서 — 교체 사유·인수인계 기간·품질 영향 평가가 없으면 나중에 R-15(선행 릴리스 결함)의 귀책을 못 가른다. (3) 하도급: 재위탁 4 \ra{} 5 승인 기록 — 개인정보 위탁·재위탁(최영주 CPO)과 연결되어 Day4 개인정보 처리 위탁 검토서의 입력이 된다. 계약 변경 로그 10건은 ⑮ 완성본. 예상 소요 5분.}
\end{frame}

\begin{frame}{5교시 핵심 정리}
\begin{enumerate}
\item 결함밀도 0.132 > 0.13 — 분모(인수 FP)와 분자(시험 결함)의 정의가 먼저, 판정은 그다음
\item 8월 평탄은 개선이 아니라 시험 정지(TC 210 \ra{} 95) — 곡선은 반드시 시험 노력과 함께 읽는다
\item 감리 지적은 등급보다 원인 귀속 칸 — 3·4번 이의 서면이 나중의 책임 귀속을 정한다
\item 이슈 11건·리스크 v1.2·잔여 EMV 7.26 < 8.0(여유 0.74 = R-10 하나) — 이내지만 경고, 사전 예고
\item 서면 회신·동시지연 분해·양방향 성과 기록 — 계약은 기록으로 관리된다
\end{enumerate}
\keymsg{6교시 실습 — 이슈 2건·리스크 재평가·감리 귀속 이의·검수 보류 회신·변경 대가를 직접 쓴다}
\note{교재 제4장·제5장 핵심 정리. 생각해 볼 질문(제5장 2 — "여러분 회사의 마지막 검수 보류는 서면으로 남아 있는가"). 6교시 준비물(워크북 §3.4·§4.4·§5.4 빈 템플릿, 감리 지적 원문 8건 사본)을 안내. 예상 소요 3분.}
\end{frame}
